\documentclass[10pt,a4paper]{amsart}
\usepackage[margin=19mm]{geometry}
\usepackage{amsmath,amssymb,mathtools}
\usepackage[scr=rsfs,cal=euler]{mathalfa}
\usepackage{xfrac}
\usepackage[unicode,hidelinks,bookmarks=false]{hyperref}
\newcommand{\ie}{i.e.\ }
\newcommand{\cf}{cf.\ }
\newcommand{\resp}{respectively\ }
\newcommand{\Subsection}{\subsection}
\providecommand{\nobreakdash}{\nobreak\hbox{-}}
\setlength{\emergencystretch}{2em}
\multlinegap0pt
\usepackage{bm}
\usepackage{bbm}
\usepackage{dsfont}
\usepackage{xspace}
\usepackage{cancel}
\usepackage{mathtools}
\usepackage{amsthm}
\makeatletter
\@ifundefined{thm}{\newtheorem{thm}{Thm}}{}
\@ifundefined{lemma}{\newtheorem{lemma}[thm]{Lemma}}{}
\makeatother
\def\T{{\mathbb T}}
\def\sT{ ^{\sharp \otimes {\rm{Tr}}}}
\title[The $K$-Theory of the Sphere with the Antipodal Involution]{The $K$-Theory of the Sphere with the Antipodal Involution}
\author{Jeffrey L Boersema}
\date{Source: arXiv:2506.16612v2 (CC BY 4.0)}
\begin{document}
\maketitle

\noindent\emph{Source and licence.} The $K$-Theory of the Sphere with the Antipodal Involution. Version arXiv:2506.16612v2 (CC BY 4.0), distributed under the Creative Commons Attribution 4.0 International licence, as recorded in the arXiv licence metadata for this exact version and shown on the abstract page https://arxiv.org/abs/2506.16612v2. Full rights notice: licenses/arxiv-2506.16612-CC-BY-4.0.txt.\par\smallskip

\noindent\textit{Editorial scope.} Counted unit: the Lemma whose source label is \texttt{transpose} in the article (source0.tex lines 970--986, source label \texttt{transpose}), delivered together with the article's own complete original proof of it (source0.tex lines 994--1037, one \texttt{proof} environment of 44 lines). No mathematics is written, repaired or re-derived: statement and proof are the article's own text line for line. The proof is detached from the statement in the source: the article states the result at lines 970--986 and prints its proof at lines 994--1037, and the proof environment's own header names the statement, so the association is the article's own. Required context retained uncounted: sharp (lines 988--990). Every definition, display and label of those lines is retained unchanged. Omitted from this package and not redistributed: 1--969, 987--987, 991--993, 1038--1313. Nothing inside a retained excerpt is omitted or altered: every retained body is delivered exactly as the article's lines are written, so no omission receipt of any kind accompanies this package. Citations: the retained excerpts carry no citation command at all, so no bibliography entry of the article is transcribed and no reference is redistributed. Reference adaptation: none was needed and none was made; every reference inside the delivered unit resolves inside it, so no unresolved reference or citation is produced. Printed slips kept literal: no printed word, symbol, hypothesis or number of the counted statement or of its proof was altered. Layout and class: the article's own class is \texttt{amsart} with packages including amsthm, amsfonts, ulem, amsmath, amssymb, mathrsfs, bbm, tikz, hhline, multicol, tabularx, multirow, inputenc, enumitem; the wrapper is the lane's pinned \texttt{amsart} editorial wrapper at 10pt on A4, which restates the theorem-like environments the retained text needs and adds the article's own macro definitions that the retained text needs (\texttt{\textbackslash T}, \texttt{\textbackslash sT}), with extra wrapper packages (bm, bbm, dsfont, xspace, cancel, mathtools). The delivered numbering is the unit's own. Assets: no retained span contains a figure environment, a graphics-inclusion command, a PDF-page inclusion, a table or a TikZ picture; no image file is copied into this package, the package declares no asset dependency and no image file is redistributed with it. Licence: version arXiv:2506.16612v2 of this article is distributed under the Creative Commons Attribution 4.0 International licence, as recorded in the arXiv licence metadata for this exact version and shown on the abstract page https://arxiv.org/abs/2506.16612v2. A full rights notice is delivered at licenses/arxiv-2506.16612-CC-BY-4.0.txt. Characters: every retained excerpt is pure ASCII, so the delivered engine, XeLaTeX with its default Latin Modern text font, reports no missing character.
\begin{lemma} \label{sharp}
The same statements in Lemma~\ref{transpose} hold if we replace $x\T$ with $x\sT$ throughout.
\end{lemma}

\begin{lemma} \label{transpose}
Let $a_1, \dots, a_k, b_1, \dots, b_\ell$ be a list of complex self-adjoint $(k + \ell)$-Clifford generators satisfying
$$(a_i)\T = a_i \quad \text{and} \quad (b_i)\T = -b_i \; .$$
Define 
$$\widetilde a_i = i (a_1 a_2 \dots \widehat a _i \dots a_k) \text{~for $1 \leq i \leq k$} 
\quad \text{and} \quad \widetilde b _i = i (b_1 b_2 \dots \widehat{b_i} \dots b_\ell) \text{~for $1 \leq i \leq \ell$} \; .
$$
Then
\begin{enumerate}
\item If $\ell \equiv 0 \pmod 4$, then
$a_1, \dots, a_k, \widetilde b_1, \dots, \widetilde b_\ell$ is a list of complex self-adjoint $(k + \ell)$-Clifford generators 
and $(\widetilde{b}_i)\T = \widetilde{b}_i$
\item If $k \equiv 0 \pmod 4$, then
$\widetilde{a}_1, \dots, \widetilde{a}_k, b_1, \dots, b_\ell$ is a list of complex self-adjoint $(k + \ell)$-Clifford generators 
and $(\widetilde{a}_i)\T = -\widetilde{a}_i$
\end{enumerate}
\end{lemma}

\begin{proof}[Proof of Lemmas~\ref{transpose} and \ref{sharp}.]
For part (1), check first that
\begin{align*}
(\widetilde b_i)^* &= -i (b_1 b_2 \dots \widehat{b_i} \dots b_\ell)^* \\
	&= -i (b_\ell \dots \widehat{b_i} \dots b_2 b_1) \\
	&= -i (-1)^{\Delta(\ell-2)} (b_1 b_2 \dots \widehat{b_i} \dots b_\ell) \\
	&= \widetilde b_i 
\end{align*}
since $\ell-2 \equiv 2 \pmod 4$. From this it follows that $\widetilde b_i$ are self-adjoint unitaries. In particular $(\widetilde b_i)^2 = 1$.

Similarly we have
\begin{align*}
(\widetilde b_i)\T &=  i (b_1 b_2 \dots \widehat{b_i} \dots b_\ell)\T \\
	& = i (-1)^{\Delta(\ell-2)} (b_1\T b_2\T \dots \widehat{b_i} \dots b_\ell\T) \\
	&= i (-1)^{\Delta(\ell-2)} (-1)^{\ell-1} (b_1 b_2 \dots \widehat{b_i} \dots b_\ell) \\
	&= \widetilde b_i \; .
\end{align*}
We also compute, for all $1 \leq i \leq k, 1 \leq j \leq \ell$
\begin{align*}
a_i \widetilde b_j &= i a_i (b_1 \dots \widehat{b}_j \dots b_\ell) \\
	&= i (-1)^{\ell-1} (b_1 \dots \widehat{b}_j \dots b_\ell) a_i \\
	&= - \widetilde b_j a_i
\end{align*}
and, for all $1 \leq i < j \leq \ell$
\begin{align*}
\widetilde b_i \widetilde b_j 
	&= i^2 (b_1 \dots \widehat{b}_i \dots b_\ell)(b_1 \dots \widehat{b}_j \dots b_\ell) \\
	&= i^2 (-1)^{\ell-1} (-1)^{\ell-2} (b_1 \dots \widehat{b}_j \dots b_\ell)(b_1 \dots \widehat{b}_i \dots b_\ell) \\
	&= - \widetilde b_j \widetilde b_i \; .
\end{align*}
This proves Part (1).

For Part (2) we set $\widetilde a_i = i (a_1 a_2 \dots \widehat a _i \dots a_k)$. 
Then we have 
\begin{align*}
(\widetilde a_i)\T &=  i (a_1 a_2 \dots \widehat{a_i} \dots a_k)\T \\
	& = i (-1)^{\Delta(k-2)} (a_1\T a_2\T \dots \widehat{a_i} \dots a_k\T) \\
	&= i (-1)^{\Delta(k-2)}  (a_1 a_2 \dots \widehat{a_i} \dots a_k) \\
	&= - \widetilde a_i; .
\end{align*}
The proof is otherwise the same  as for Part (1). 

The proof of Lemma~\ref{sharp} is identical to that of Lemma~\ref{transpose}, since both involutions ${\rm Tr}$ and $\sharp \otimes {\rm Tr}$ have the same formal properties -- they are linear and antimultiplicative.
\end{proof}

\end{document}
